\documentclass[11pt,a4paper]{article}
\usepackage[margin=2.8cm]{geometry}
\usepackage{amsmath,amssymb,amsthm,mathrsfs,mathtools}
\usepackage[T1]{fontenc}
\usepackage{microtype}
\usepackage{booktabs}
\usepackage[colorlinks=true,linkcolor=blue,citecolor=blue,urlcolor=blue,hypertexnames=false]{hyperref}

\newtheorem{theorem}{Theorem}[section]
\newtheorem{proposition}[theorem]{Proposition}
\newtheorem{lemma}[theorem]{Lemma}
\newtheorem{corollary}[theorem]{Corollary}
\newtheorem{hypothesis}[theorem]{Hypothesis}
\theoremstyle{remark}
\newtheorem{remark}[theorem]{Remark}

\newcommand{\R}{\mathbb R}
\newcommand{\A}{\mathfrak A}
\newcommand{\Hh}{\mathcal H_0}
\newcommand{\Om}{\Omega}
\newcommand{\G}{\mathcal G}
\newcommand{\Q}{\mathcal Q}
\newcommand{\E}{\mathbb E}
\DeclareMathOperator{\supp}{supp}
\DeclareMathOperator{\spec}{spec}
\DeclareMathOperator{\dist}{dist}
\DeclareMathOperator{\Var}{Var}
\DeclareMathOperator{\Tr}{Tr}

\title{M4b: discharge audit and corrected proof program\\
for continuum LPPL in $P(\varphi)_2$}
\author{}
\date{8 August 2026}

\begin{document}
\maketitle

\begin{abstract}
This note audits, one by one, the four inputs left open by the conditional LPPL theorem in
\texttt{lppl.tex}.  Two conclusions are unconditional within the standard spatial-cutoff
$P(\varphi)_2$ construction.  First, the actual cutoff dynamics has propagation speed at most one;
finite propagation is therefore a theorem, not a hypothesis.  Second, every cutoff ground-state
vector belongs to the operator domain of every compactly smeared Wick polynomial.  This removes
the vector-domain part of the differentiability input and the domain clause from the moment input.

The other conclusions are negative but precise.  A common quadratic-form domain on the entire
closed interpolation interval is not a legitimate generic assumption when an interaction is
turned on from zero.  It is replaced below by local form differentiability on $0<s<1$ and endpoint
norm-resolvent continuity.  The proposed modern log-Sobolev input has the unweighted stochastic-
quantisation Dirichlet form, not the covariance-weighted form used in \texttt{h6-attack.tex}; it
therefore proves neither the physical Hamiltonian gap nor the required local second moment.
Moreover, the induction in \texttt{h6-attack.tex} contains two false operator estimates and is
withdrawn.  A corrected conditional induction is given, exposing the additional one-point bound
that would actually be needed.

Thus the continuum LPPL argument itself is complete, exact propagation is complete, and the
ground-state domain question is complete.  The affine path regularity is also proved below from
the Feynman--Kac--Nelson formula once the same uniform gap already assumed by LPPL is available.
Two sharply stated analytic obligations remain: a uniform physical gap in a specified noncritical
regime and a uniform local second moment.  Neither is asserted without proof.
\end{abstract}

\section{Setting and verdict}

Let
\[
 K(g)=H_0+V(g),\qquad
 V(h)=\int_{\R}h(x):P(\varphi_0(x)):\,dx,
 \qquad H(g)=K(g)-E(g),
\]
where $E(g)=\inf\spec K(g)$ and
\[
 \G=\{g\in C_c^\infty(\R):0\le g\le1,\quad
             g=1\text{ on a nonempty bounded interval}\}.
\]
The ground state of $H(g)$ is denoted by $\Om_g$ and is normalized.  The time-zero local algebra
over a bounded open interval $O$ is $\A(O)$, and
\[
 \beta_t^g(A)=e^{itH(g)}Ae^{-itH(g)}.
\]
All assertions below use the standard $Q$-space construction of the spatially cutoff
$P(\varphi)_2$ Hamiltonian: $V(h)$ is the self-adjoint multiplication operator associated with the
Wick polynomial, and the usual $L^1\cap L^2$ extension of the cutoff is allowed when a
characteristic function is used in a localization argument.

\begin{center}
\begin{tabular}{@{}lll@{}}
\toprule
Input in \texttt{lppl.tex} & Verdict & Exact remainder\\
\midrule
finite propagation (L) & proved, Proposition~\ref{prop:propagation} & none\\
$\Om_g\in D(V(h))$ & proved, Proposition~\ref{prop:domain} & none\\
path differentiability (D) & follows from (G) & no independent input, \S\ref{sec:path}\\
uniform physical gap (G) & open & (G$_\mathrm{exact}$), \S\ref{sec:gap}\\
uniform local moment (M) & open & (M$_2$), \S\ref{sec:moment}\\
\bottomrule
\end{tabular}
\end{center}

The word ``open'' in the table is mathematical, not procedural: the sources audited here do not
contain the required gap or moment theorem, and the arguments previously proposed in this folder
do not prove either one.  The conditional LPPL theorem remains valid after finite propagation and
the vector-domain condition are removed from its hypothesis list; Theorem~\ref{thm:FKNpath} also
removes path differentiability as an independent hypothesis.

\section{Exact finite propagation for the cutoff dynamics}

We spell out the localization argument because the LPPL cell expansion uses the cutoff dynamics
far outside the plateau $\{g=1\}$.

\begin{proposition}[cutoff propagation]\label{prop:propagation}
For every $g\in\G$, every bounded open interval $O\subset\R$, and every $t\in\R$,
\[
 \beta_t^g\bigl(\A(O)\bigr)\subseteq\A(O_{|t|}),
 \qquad O_r=\{x:\dist(x,O)<r\}.
\]
The speed bound is one and is independent of $g$.  No assumption that $g=1$ on $O_{|t|}$ is
needed for this inclusion.
\end{proposition}

\begin{proof}
The scalar $E(g)I$ cancels from the adjoint action, so it suffices to use $K(g)=H_0+V(g)$.
Write
\[
 \beta_t^0=\operatorname{Ad}e^{itH_0},\qquad
 \beta_t^{I,g}=\operatorname{Ad}e^{itV(g)}.
\]

First, $\beta^{I,g}$ has propagation speed zero.  It is enough initially to take an observable
$A$ generated by Weyl operators with test functions supported in a compact set $B\Subset O$.
Choose $\chi\in C_c^\infty(O)$ with $\chi=1$ on a neighbourhood of $B$ and split
\[
 V(g)=V(\chi g)+V((1-\chi)g).
\]
The two summands strongly commute because they are multiplication operators in $Q$-space.
Affiliation gives
\[
 e^{itV(\chi g)}\in\A(O),
 \qquad e^{itV((1-\chi)g)}\in\A(B)'.
\]
Consequently
\[
 \beta_t^{I,g}(A)
 =e^{itV(\chi g)}A e^{-itV(\chi g)}\in\A(O).
\]
Strong density of the compactly supported Weyl-generated subalgebra and strong closedness of
$\A(O)$ extend this conclusion to every $A\in\A(O)$.

Second, the free Klein--Gordon equation has propagation speed one, hence
\[
 \beta_t^0(\A(O))\subseteq\A(O_{|t|}).
\]
For $A\in\A(O)$, the $n$th Trotter approximant
\[
 A_n(t)=\bigl(\beta_{t/n}^0\circ\beta_{t/n}^{I,g}\bigr)^n(A)
\]
therefore belongs to $\A(O_{|t|})$: each interaction step leaves the current localization region
unchanged, and each free step enlarges it by $|t|/n$.  The self-adjointness theorem for
$H_0+V(g)$ and the Trotter product formula give
\[
 e^{itK(g)}=\mathop{\mathrm{s\! -\! lim}}_{n\to\infty}
             \bigl(e^{itH_0/n}e^{itV(g)/n}\bigr)^n.
\]
The same holds for the adjoints.  Hence $A_n(t)\to\beta_t^g(A)$ strongly.  Since
$\A(O_{|t|})$ is a von Neumann algebra and therefore strongly closed,
$\beta_t^g(A)\in\A(O_{|t|})$.
\end{proof}

\begin{remark}
This is the argument of Glimm--Jaffe's finite-propagation theorem
\cite[Ch.~4, Prop.~4.1.1 and Thm.~4.1.2]{GJExp}; their later formulation
\cite[Thm.~8.1]{GJModels} also proves cutoff independence when $g=1$ on the relevant light cone.
Cutoff independence is an additional conclusion.  The propagation inclusion itself follows from
``free speed one + interaction speed zero'' for every spatial cutoff $g$.
\end{remark}

\section{The ground state is in every local Wick-polynomial domain}

\begin{proposition}[domain smoothing]\label{prop:domain}
Let $g\in\G$ and let $h\in C_c^\infty(\R)$ be real.  Then
\[
 \Om_g\in D(V(h)).
\]
In particular, for an affine path $g_s=(1-s)g+sg'$ and $v=g'-g$,
$\Om_{g_s}\in D(V(v))$ for every $0\le s\le1$.
\end{proposition}

\begin{proof}
The self-adjointness and smoothing theorem for cutoff $P(\varphi)_2$ Hamiltonians states that, if
$K=H_0+V(g)$ and $V'$ is any compactly smeared Wick polynomial in the standard polynomial
$L^2(Q)$ class, then for every $t>0$,
\[
 \operatorname{Ran}e^{-tK}\subseteq D(H_0)\cap D(V').
\]
See \cite[Ch.~7, Thm.~7.3 and Lem.~7.4]{GJModels}.  Take $V'=V(h)$.  Since
$K(g)\Om_g=E(g)\Om_g$,
\[
 \Om_g=e^{tE(g)}e^{-tK(g)}\Om_g.
\]
The right-hand side belongs to $D(V(h))$ by the displayed range inclusion.  This proves the
claim.  Taking $h=v$ gives the path statement.
\end{proof}

Thus the domain clause in hypothesis (M) is redundant.  The unresolved content of (M) is solely
the numerical estimate $\sup_g\|V(h)\Om_g\|<\infty$ for unit-cell $h$.

\section{The correct differentiability input}\label{sec:path}

\subsection{Why a global common form domain is not available automatically}

The original hypothesis (D) in \texttt{lppl.tex} imposed one common form domain for all
$0\le s\le1$.  Turning on a positive highest-degree interaction can change the form domain at an
endpoint.  This is already visible in one-dimensional quantum mechanics.

\begin{proposition}[explicit endpoint obstruction]\label{prop:oscillator}
On $L^2(\R)$ let
\[
 K_s=-\frac{d^2}{dx^2}+x^2+s x^4,\qquad 0\le s\le1,
\]
with the Friedrichs realization.  Its form domains satisfy
\[
 D(q_0)=\{\psi\in H^1(\R):x\psi\in L^2\},
\]
whereas, for every $s>0$,
\[
 D(q_s)=\{\psi\in H^1(\R):x\psi\in L^2,\ x^2\psi\in L^2\}.
\]
In particular the domains are not common on $[0,1]$.
\end{proposition}

\begin{proof}
The two identities follow directly from the nonnegative closed forms
\[
 q_s[\psi]=\|\psi'\|_2^2+\|x\psi\|_2^2+s\|x^2\psi\|_2^2.
\]
To see that the inclusion is strict, put $\psi(x)=(1+x^2)^{-1}$.  Then
$\psi,\psi',x\psi\in L^2(\R)$: at infinity their absolute values are respectively
$O(|x|^{-2})$, $O(|x|^{-3})$, and $O(|x|^{-1})$.  But
\[
 |x^2\psi(x)|^2=\frac{x^4}{(1+x^2)^2}\longrightarrow1,
\]
so $x^2\psi\notin L^2(\R)$.  Hence $\psi\in D(q_0)\setminus D(q_s)$ for $s>0$.
\end{proof}

This does not show that the $P(\varphi)_2$ ground state is discontinuous.  It shows only that a
proof may not assume a global common domain without a separate theorem.

\subsection{A minimal replacement}

Let $g,g'\in\G$ have a common plateau, $g_s=(1-s)g+sg'$, $v=g'-g$, and write
$K_s=K(g_s)$, $E_s=E(g_s)$, $H_s=K_s-E_s$, and
$P_s=|\Om_s\rangle\langle\Om_s|$.

\begin{hypothesis}[minimal path regularity, D$_\mathrm{min}$]\label{hyp:Dmin}
The following two statements hold.
\begin{enumerate}
\item[(D1)] For every $s_0\in(0,1)$ there are an open interval $I\ni s_0$ and a Hilbert form
domain $\Q_I\hookrightarrow\Hh$, common to $s\in I$, such that the closed forms of $K_s$ obey
\[
 k_s=k_{s_0}+(s-s_0)v_I,
\]
where $v_I$ is a bounded symmetric form on $\Q_I$.  On the ground state it is represented by the
operator $V(v)$:
\[
 v_I(\xi,\Om_s)=\langle\xi,V(v)\Om_s\rangle
 \quad(\xi\in\Q_I,\ s\in I).
\]
\item[(D2)] The shifted Hamiltonians converge in the strong resolvent sense at both endpoints:
\[
 H_s\xrightarrow[s\downarrow0]{\rm s.r.}H_0,
 \qquad H_s\xrightarrow[s\uparrow1]{\rm s.r.}H_1.
\]
\end{enumerate}
\end{hypothesis}

The notation $H_0$ in (D2) means $H(g)$, not the free Hamiltonian.  Proposition~\ref{prop:domain}
already supplies the operator-domain assertion in (D1).

\begin{lemma}[local form regularity gives the derivative]\label{lem:localderivative}
Assume (D1) and suppose
\[
 \spec H_s\subseteq\{0\}\cup[\gamma,\infty)
 \qquad(s\in I)
\]
for some $\gamma>0$.  After a phase choice, $s\mapsto\Om_s$ is $C^1$ locally on $I$ and
\begin{align}
 \dot E_s&=\langle\Om_s,V(v)\Om_s\rangle,\label{eq:HF}\tag{4.1}\\
 \dot\Om_s&=-H_s^{-1}(1-P_s)V(v)\Om_s.\label{eq:vector}\tag{4.2}
\end{align}
\end{lemma}

\begin{proof}
Give $\Q_I$ the norm
$\|\xi\|_{\Q_I}^2=k_{s_0}[\xi]+c\|\xi\|^2$ with $c$ so large that it is positive.  Let
$J:\Q_I\to\Q_I^*$ be the canonical embedding induced by the Hilbert-space inner product, and
let $A_s:\Q_I\to\Q_I^*$ represent $k_s$.  Then
\[
 A_s=A_{s_0}+(s-s_0)V_I
\]
in $B(\Q_I,\Q_I^*)$.  If $z$ is in the resolvent set of $K_s$, the weak resolvent is
$(A_s-zJ)^{-1}:\Q_I^*\to\Q_I$.  The inverse identity gives, whenever both sides exist,
\[
 (A_{s+h}-zJ)^{-1}-(A_s-zJ)^{-1}
 =-h(A_{s+h}-zJ)^{-1}V_I(A_s-zJ)^{-1}.
\]
Thus the resolvent is $C^1$ as a map $\Q_I^*\to\Q_I$, with derivative
\[
 - (A_s-zJ)^{-1}V_I(A_s-zJ)^{-1}.
\]
Composing with the continuous embeddings $\Hh\hookrightarrow\Q_I^*$ and
$\Q_I\hookrightarrow\Hh$ gives norm-resolvent $C^1$ dependence on $s$.

Fix $s_0$ and choose a circle $\Gamma$ enclosing only the simple eigenvalue $E_s$ for all $s$ in
a smaller neighbourhood; the gap and norm-resolvent continuity permit this.  The Riesz projection
\[
 P_s=-\frac{1}{2\pi i}\int_\Gamma(K_s-z)^{-1}\,dz
\]
is $C^1$ in $B(\Hh)$ and, by the preceding mapping statement, as a map $\Hh\to\Q_I$.  The vector
\[
 \widetilde\Om_s=\frac{P_s\Om_{s_0}}{\|P_s\Om_{s_0}\|}
\]
is therefore $C^1$ in $\Q_I$ after shrinking the neighbourhood.  Multiplication by a $C^1$ phase
puts it in the parallel gauge $\langle\Om_s,\dot\Om_s\rangle=0$.

Differentiate the weak eigenvalue equation
\[
 k_s(\xi,\Om_s)=E_s\langle\xi,\Om_s\rangle,
 \qquad \xi\in\Q_I.
\]
Taking $\xi=\Om_s$ after differentiating and using the weak eigenvalue equation once more with
test vector $\dot\Om_s$ gives \eqref{eq:HF}.  For general $\xi$ one obtains
\[
 (k_s-E_s\langle\cdot,\cdot\rangle)(\xi,\dot\Om_s)
 =-\langle\xi,(V(v)-\dot E_s)\Om_s\rangle.
\]
The right side is represented by a Hilbert-space vector by Proposition~\ref{prop:domain}.
The representation theorem for closed forms therefore gives
\[
 H_s\dot\Om_s=-(1-P_s)V(v)\Om_s.
\]
The parallel gauge puts $\dot\Om_s$ in $(1-P_s)\Hh$, where $H_s^{-1}$ has norm at most
$\gamma^{-1}$.  This proves \eqref{eq:vector}.
\end{proof}

\begin{lemma}[endpoint projection continuity]\label{lem:endpoint}
Assume (D2) and a uniform gap
$\spec H_s\subseteq\{0\}\cup[\gamma,\infty)$ for $0\le s\le1$.  Then
\[
 \|P_s-P_0\|\longrightarrow0\quad(s\downarrow0),
 \qquad
 \|P_s-P_1\|\longrightarrow0\quad(s\uparrow1).
\]
Consequently $s\mapsto\omega_{g_s}(A)$ is continuous at both endpoints for every bounded $A$.
\end{lemma}

\begin{proof}
Strong resolvent convergence implies
$e^{-tH_s}\to e^{-tH_0}$ strongly for every fixed $t>0$.  The uniform gap gives
\[
 \|e^{-tH_s}-P_s\|\le e^{-\gamma t},
 \qquad \|e^{-tH_0}-P_0\|\le e^{-\gamma t}.
\]
For every $\xi\in\Hh$,
\[
 \|(P_s-P_0)\xi\|
 \le2e^{-\gamma t}\|\xi\|
   +\|(e^{-tH_s}-e^{-tH_0})\xi\|.
\]
First let $s\downarrow0$, then $t\to\infty$; this proves $P_s\to P_0$ strongly.  Since both are
rank-one projections,
\[
 \|P_s-P_0\|^2=1-\Tr(P_sP_0).
\]
Strong convergence applied to a unit vector spanning $\operatorname{Ran}P_0$ makes the right-hand
side tend to zero.  The proof at $s=1$ is identical.  Finally
$\omega_{g_s}(A)=\Tr(P_sA)$ and
$|\Tr((P_s-P_0)A)|\le\|P_s-P_0\|_1\|A\|=2\|P_s-P_0\|\|A\|$ for rank at most two.
\end{proof}

\begin{corollary}[what LPPL actually needs]\label{cor:Dmin}
In \texttt{lppl.tex}, global hypothesis (D) may be replaced by
(D$_\mathrm{min}$).  The derivative formula holds for $0<s<1$.  If its right-hand side is bounded
in absolute value by a constant $C_A$ independent of $s$, then
\[
 |\omega_{g'}(A)-\omega_g(A)|\le C_A.
\]
\end{corollary}

\begin{proof}
Lemmas~\ref{lem:localderivative} and \ref{lem:endpoint} give differentiability on $(0,1)$ and
continuity on $[0,1]$.  For $0<\varepsilon<1/2$, the ordinary mean-value estimate on
$[\varepsilon,1-\varepsilon]$ gives
\[
 |\omega_{g_{1-\varepsilon}}(A)-\omega_{g_\varepsilon}(A)|
 \le (1-2\varepsilon)C_A.
\]
Let $\varepsilon\downarrow0$.
\end{proof}

The form criterion above is sufficient, but the actual model admits a cleaner route which bypasses
form domains altogether.

\begin{theorem}[FKN path regularity]\label{thm:FKNpath}
Let $g,g'\in\G$ have a common nonempty plateau, let $g_s=(1-s)g+sg'$, and put
$v=g'-g$.  Assume the uniform gap
\[
 \spec H(g_s)\subseteq\{0\}\cup[\gamma,\infty)
 \qquad(0\le s\le1).
\]
Then the ground-state projections $P_s$ are norm-continuous on $[0,1]$ and $C^1$ on $(0,1)$.
Locally on $(0,1)$, normalized vectors can be chosen $C^1$ in Hilbert norm and in the parallel
gauge, and they satisfy
\begin{align*}
 \dot E_s&=\langle\Om_s,V(v)\Om_s\rangle,\\
 \dot\Om_s&=-H_s^{-1}(1-P_s)V(v)\Om_s.
\end{align*}
Moreover $H_s\to H_0$ in strong resolvent sense as $s\downarrow0$ and
$H_s\to H_1$ in strong resolvent sense as $s\uparrow1$.  Thus affine-path differentiability is
not an independent hypothesis once the LPPL gap hypothesis is assumed.
\end{theorem}

\begin{proof}
Fix a Euclidean time $T>0$.  On free Euclidean path space let $J_0,J_T$ be the time-slice maps and
put
\[
 \mathcal U_s=\int_0^T\!\int_{\R}g_s(x):P(\varphi(t,x)):\,dx\,dt,
 \qquad w_s=e^{-\mathcal U_s},
 \qquad \mathcal U_v=\mathcal U_1-\mathcal U_0.
\]
The Feynman--Kac--Nelson formula gives, initially for bounded time-zero functions $f,h$,
\begin{equation}\label{eq:FKNpath}
 \langle f,e^{-TK(g_s)}h\rangle
 =\int \overline{J_Tf}\,(J_0h)w_s\,d\mu_0.                    \tag{4.3}
\end{equation}
Let
\[
 q_T=\frac{2}{1-e^{-m_0T}},\qquad r_T=\frac{q_T}{q_T-1}.
\]
For completeness, the time-slice product estimate follows with no hidden interpolation step.
Let $\nu_0$ be the time-zero Gaussian marginal of $\mu_0$, let
$\mathsf P_T=e^{-TH_0}$ be its Markov semigroup, and set
\[
 r_T=\frac{2}{1+e^{-m_0T}},\qquad
 p_T=\frac2{r_T}=1+e^{-m_0T},\qquad
 p_T'=1+e^{m_0T}.
\]
Then $p_T'$ is conjugate to $p_T$ and
$(p_T'-1)/(p_T-1)=e^{2m_0T}$.  Nelson hypercontractivity therefore gives
$\|\mathsf P_TF\|_{p_T'}\le\|F\|_{p_T}$.  The Markov property and H\"older's inequality now
give, first for bounded $f,h$,
\begin{align*}
 \|(J_Tf)(J_0h)\|_{r_T}^{r_T}
 &=\int |h|^{r_T}\,\mathsf P_T(|f|^{r_T})\,d\nu_0\\
 &\le \||h|^{r_T}\|_{p_T}
       \|\mathsf P_T(|f|^{r_T})\|_{p_T'}\\
 &\le \||h|^{r_T}\|_{p_T}\||f|^{r_T}\|_{p_T}
 =\|h\|_2^{r_T}\|f\|_2^{r_T}.
\end{align*}
Taking the $r_T$th root proves
\begin{equation}\label{eq:sliceproduct}
 \|(J_Tf)(J_0h)\|_{r_T}\le\|f\|_2\|h\|_2.                    \tag{4.4}
\end{equation}
This is the estimate used in the transfer-matrix construction
\cite[Thms.~V.10 and V.12]{Simon}.

Every compactly smeared Wick polynomial belongs to $L^p(\mu_0)$ for every finite $p$; hence
$\mathcal U_v\in\bigcap_{p<\infty}L^p(\mu_0)$.  Since $g_0,g_1\ge0$ and $P$ is semibounded, the
standard exponential estimate \cite[Thm.~V.7]{Simon} gives
$e^{-a\mathcal U_0},e^{-a\mathcal U_1}\in L^1(\mu_0)$ for every $a>0$.  Linearity
$\mathcal U_s=(1-s)\mathcal U_0+s\mathcal U_1$ and H\"older's inequality therefore imply
\begin{equation}\label{eq:weightbound}
 \|w_s\|_a^a
 \le\left(\int e^{-a\mathcal U_0}d\mu_0\right)^{1-s}
     \left(\int e^{-a\mathcal U_1}d\mu_0\right)^s             \tag{4.5}
\end{equation}
for every finite $a$.  In particular the right side is bounded uniformly in $0\le s\le1$.

Because $\mathcal U_s-\mathcal U_r=(s-r)\mathcal U_v$ in every $L^p$, $w_s\to w_r$ in
probability as $s\to r$.  Applying \eqref{eq:weightbound} with an exponent strictly larger than
$a$ gives uniform integrability of $|w_s|^a$; Vitali's theorem yields
\begin{equation}\label{eq:weightcont}
 \|w_s-w_r\|_a\longrightarrow0\qquad(s\to r)                 \tag{4.6}
\end{equation}
for every finite $a$.

Fix $q=q_T$ and choose finite $p,a$ with $q^{-1}=p^{-1}+a^{-1}$.  Pointwise in path space,
\[
 \frac{w_{s+h}-w_s}{h}
 =-\mathcal U_v\int_0^1w_{s+\theta h}\,d\theta.
\]
Using H\"older's inequality and \eqref{eq:weightcont},
\begin{align*}
 \left\|\frac{w_{s+h}-w_s}{h}+\mathcal U_vw_s\right\|_q
 &\le\|\mathcal U_v\|_p
       \sup_{0\le\theta\le1}\|w_{s+\theta h}-w_s\|_a
 \longrightarrow0.
\end{align*}
Thus $s\mapsto w_s$ is $C^1$ in $L^{q_T}$ on $(0,1)$, with one-sided continuity at the endpoints,
and $\dot w_s=-\mathcal U_vw_s$.  Indeed, continuity of the derivative follows from
\[
 \|\mathcal U_v(w_s-w_r)\|_{q_T}
 \le\|\mathcal U_v\|_p\|w_s-w_r\|_a\longrightarrow0.
\]
Equations \eqref{eq:FKNpath}--\eqref{eq:sliceproduct} now give
\[
 \|e^{-TK(g_s)}-e^{-TK(g_r)}\|
 \le\|w_s-w_r\|_{q_T}
\]
and the same estimate for difference quotients.  Therefore
\[
 s\longmapsto e^{-TK(g_s)}
\]
is norm-continuous on $[0,1]$ and $C^1$ on $(0,1)$.

Set $L_s=e^{-TK(g_s)}$.  Its largest spectral value is
$\lambda_s=\|L_s\|=e^{-TE_s}$, it is simple, and the gap assumption puts the remainder of its
spectrum in $[0,e^{-T\gamma}\lambda_s]$.  Norm continuity of $L_s$ implies continuity of
$\lambda_s$, hence of $E_s$.  A local Riesz contour around $\lambda_s$ then shows that $P_s$ is
$C^1$ on $(0,1)$ and norm-continuous at both endpoints.  A local vector is obtained from
$P_s\Om_{s_0}/\|P_s\Om_{s_0}\|$; a $C^1$ phase change imposes
$\langle\Om_s,\dot\Om_s\rangle=0$.  Since
$\lambda_s=\langle\Om_s,L_s\Om_s\rangle$, the eigenvalue $\lambda_s$ and therefore
$E_s=-T^{-1}\log\lambda_s$ are $C^1$ on $(0,1)$.

It remains to identify the derivative.  The cutoff Hamiltonians have the common operator core
$\mathscr C=C^\infty(H_0)$ and, on this core,
\[
 K(g_s)\xi=K(g_{s_0})\xi+(s-s_0)V(v)\xi
 \qquad(\xi\in\mathscr C),
\]
see \cite[Thm.~V.12(a)]{Simon}.  For $\xi\in\mathscr C$, differentiate
\[
 \langle K(g_s)\xi,\Om_s\rangle=E_s\langle\xi,\Om_s\rangle.
\]
Proposition~\ref{prop:domain} permits
$\langle V(v)\xi,\Om_s\rangle=\langle\xi,V(v)\Om_s\rangle$, so
\[
 \langle (K(g_s)-E_s)\xi,\dot\Om_s\rangle
 =\langle\xi,(\dot E_s-V(v))\Om_s\rangle.
\]
Because $\mathscr C$ is an operator core, the definition of the adjoint gives
$\dot\Om_s\in D(H_s)$ and
\[
 H_s\dot\Om_s=(\dot E_s-V(v))\Om_s.
\]
Taking the inner product with $\Om_s$ proves the Hellmann--Feynman identity.  The parallel gauge
and the gap then give the asserted formula for $\dot\Om_s$.

Finally, the preceding norm-continuity argument holds for every $T>0$.  Since $E_s$ is continuous,
$e^{-TH_s}=e^{TE_s}e^{-TK(g_s)}$ converges in norm at the endpoints for every $T>0$.  Semigroup
convergence gives the explicit resolvent identity
\[
 (H_s+1)^{-1}=\int_0^\infty e^{-T}e^{-TH_s}\,dT.
\]
The integrand differences converge in norm pointwise in $T>0$ and are bounded by $2e^{-T}$.
Bochner dominated convergence therefore proves norm-resolvent, hence strong-resolvent,
convergence at both endpoints.
\end{proof}

\section{The uniform-gap obligation}\label{sec:gap}

\subsection{What the cited log-Sobolev theorem says}

Bauerschmidt--Dagallier \cite[Eq.~(1.6), Thm.~1.1]{BD} use, on a lattice of mesh $\varepsilon$,
the Dirichlet form
\[
 \mathcal D_{\rm BD}(F)
 =\varepsilon^{-d}\sum_x
   \E\!\left[\left(\frac{\partial F}{\partial\varphi_x}\right)^2\right].
\]
It is the unweighted $\ell^2$ gradient form of stochastic quantisation.  The covariance-weighted
form used in \texttt{h6-attack.tex} is instead
\[
 \mathcal D_C(F)=\E\langle\nabla F,C\nabla F\rangle,
 \qquad \widehat C(k)=\frac{1}{2\sqrt{k^2+m_0^2}}.
\]
Since $0<C\le(2m_0)^{-1}I$,
\[
 \mathcal D_C(F)\le(2m_0)^{-1}\mathcal D_{L^2}(F).
\]
Thus a covariance-weighted Poincar\'e inequality implies an unweighted one (with a changed
constant), but the converse does not follow: $\widehat C(k)\to0$ as $|k|\to\infty$, so no
inequality $I\le cC$ is possible.  The mismatch is structural, not a missing constant.

Furthermore, the spectral gap furnished by an LSI is the gap of the Markov generator associated
with $\mathcal D_{\rm BD}$.  That generator is not the physical Hamiltonian $H(g)$.  A separate
Euclidean-time correlation argument is required.

\subsection{The exact Euclidean-to-Hamiltonian implication}

\begin{lemma}[dense-family decay forces a Hamiltonian gap]\label{lem:decaygap}
Let $H\ge0$ have a simple ground state $\Om$, let $Q=1-|\Om\rangle\langle\Om|$, and let
$\mathscr D\subset Q\mathcal H$ be dense.  Suppose there is $\gamma>0$ such that for every
$\psi\in\mathscr D$ there is $C_\psi<\infty$ satisfying
\[
 0\le\langle\psi,e^{-tH}\psi\rangle\le C_\psi e^{-\gamma t}
 \qquad(t\ge0).
\]
Then $\spec(H|_{Q\mathcal H})\subseteq[\gamma,\infty)$.
\end{lemma}

\begin{proof}
Fix $0<\varepsilon<\gamma$ and let
$R_\varepsilon=\mathbf1_{[0,\gamma-\varepsilon]}(H)Q$.  For $\psi\in\mathscr D$, the spectral
theorem gives
\[
 \langle\psi,e^{-tH}\psi\rangle
 \ge e^{-(\gamma-\varepsilon)t}\|R_\varepsilon\psi\|^2.
\]
Combining this with the assumed upper bound yields
\[
 \|R_\varepsilon\psi\|^2\le C_\psi e^{-\varepsilon t}
 \quad(t\ge0),
\]
and therefore $R_\varepsilon\psi=0$.  Density of $\mathscr D$ and boundedness of
$R_\varepsilon$ imply $R_\varepsilon=0$.  Since this holds for every $\varepsilon>0$, the spectral
projection of $H|_{Q\mathcal H}$ on $[0,\gamma)$ is zero.
\end{proof}

Decay of the two-point function of the single field $\varphi$ is not enough to invoke this lemma:
it controls only vectors of the form $\varphi(f)\Om$.  Cyclicity of the full local algebra does not
turn that one-particle family into a dense family.  One needs the same exponent for a collection
of centered polynomial/cylinder observables whose vectors are dense.

The unresolved gap statement is consequently:

\begin{quote}
\textbf{(G$_\mathrm{exact}$).} Specify a noncritical class of polynomials/couplings, a cylinder
$*$-algebra $\mathscr C$, and $\gamma>0$ such that, for every $g\in\G$,
\[
 \overline{\operatorname{span}}\{F\Om_g:F\in\mathscr C\}=\Hh
\]
and, for every $F\in\mathscr C$, there is $C_{g,F}<\infty$ for which
\[
 0\le
 \langle Q_gF\Om_g,e^{-tH(g)}Q_gF\Om_g\rangle
 \le C_{g,F}e^{-\gamma t}
 \qquad(t\ge0),
 \quad Q_g=1-|\Om_g\rangle\langle\Om_g|.
\]
Only the exponent $\gamma$ must be uniform in $g$ and $F$.  Then
$\operatorname{span}\{Q_gF\Om_g:F\in\mathscr C\}$ is dense in $Q_g\Hh$, and
Lemma~\ref{lem:decaygap} gives
$\spec H(g)\subseteq\{0\}\cup[\gamma,\infty)$ uniformly in $g$.
\end{quote}

The theorem of \cite{BD} is uniform in its lattice and volume parameters under its hypotheses,
but it treats the unweighted stochastic-quantisation form and constant-coupling finite-volume
measures.  It does not state (G$_\mathrm{exact}$) for arbitrary profiles $0\le g(x)\le1$.
In a symmetry-broken regime a uniform full-net gap is in any event impossible by the tunnelling
argument already proved in \texttt{../m3/m3.tex}.  Therefore the regime must be part of the gap
theorem, not appended after it.

\section{The local-second-moment obligation}\label{sec:moment}

The remaining moment statement is
\begin{equation}\label{eq:M2}
 \tag{M$_2$}
 \sup_{\substack{g\in\G,\ h\in C_c^\infty(\R;\R)\\
                   \|h\|_\infty\le1,
                   \ \supp h\subset I_h,\ |I_h|\le2}}
 \|V(h)\Om_g\|^2<\infty
\end{equation}
with one constant for all indicated $g$ and $h$.
Proposition~\ref{prop:domain} proves that every term in \eqref{eq:M2} is defined; it does not prove
uniform boundedness.

\subsection{Why the claimed Poincar\'e induction is invalid}

Let $B$ be a bounded interval, let
$C=(2\sqrt{-\partial_x^2+m_0^2})^{-1}$, and let
\[
 W_u=M_uCM_u
\]
on $L^2(B)$, where $u\in C_c^\infty(B)$ is nonzero.  The proof of Theorem~2.3 in
\texttt{h6-attack.tex} uses both
\[
 \|v\|_\infty^2\le\kappa_B\|v\|_2^2
 \quad\text{for all }v\in L^2(B),                                      \tag{6.1}
\]
and $\Tr W_u<\infty$.  Both assertions are false.

\begin{proposition}[the two failed estimates]\label{prop:failed}
There is no finite $\kappa_B$ for which \emph{(6.1)} holds.  Moreover $W_u$ is positive and
Hilbert--Schmidt but not trace class; in particular $\Tr W_u=+\infty$.
\end{proposition}

\begin{proof}
Choose $x_0\in B$ and a nonzero $\eta\in C_c^\infty((-1,1))$.  For all sufficiently large $n$,
\[
 v_n(x)=n^{1/2}\eta(n(x-x_0))
\]
is supported in $B$, while
\[
 \|v_n\|_2=\|\eta\|_2,
 \qquad \|v_n\|_\infty=n^{1/2}\|\eta\|_\infty.
\]
This contradicts (6.1).

The kernel $C(x-y)$ has only a logarithmic diagonal singularity, so
$u(x)C(x-y)u(y)\in L^2(B\times B)$; hence $W_u$ is Hilbert--Schmidt.  To compute its trace, let
$C_\Lambda$ be the positive Fourier multiplier
\[
 \widehat C_\Lambda(k)
 =\mathbf1_{[-\Lambda,\Lambda]}(k)\frac1{2\sqrt{k^2+m_0^2}}.
\]
Then $0\le M_uC_\Lambda M_u\le W_u$.  More explicitly,
$A_\Lambda=C_\Lambda^{1/2}M_u$ is Hilbert--Schmidt because
\[
 \|A_\Lambda\|_2^2
 =\frac{\|u\|_2^2}{2\pi}
   \int_{-\Lambda}^{\Lambda}\frac{dk}{2\sqrt{k^2+m_0^2}}<\infty.
\]
Hence $M_uC_\Lambda M_u=A_\Lambda^*A_\Lambda$ is positive trace class and
\begin{align*}
 \Tr(M_uC_\Lambda M_u)
 &=C_\Lambda(0)\int_Bu(x)^2\,dx\\
 &=\frac{\|u\|_2^2}{2\pi}
   \int_{-\Lambda}^{\Lambda}\frac{dk}{2\sqrt{k^2+m_0^2}}.
\end{align*}
The last integral diverges logarithmically as $\Lambda\to\infty$.  If $W_u$ were trace class,
monotonicity of the trace on positive operators would bound all these traces by $\Tr W_u$, a
contradiction.
\end{proof}

The invalid step in \texttt{h6-attack.tex} is precisely
\[
 \sum_\alpha\lambda_\alpha\|e_\alpha\|_\infty^2
 \lesssim\sum_\alpha\lambda_\alpha=\Tr W_u.
\]
Proposition~\ref{prop:failed} shows that neither inequality is available.  In addition, the
spectral vectors $e_\alpha$ are a priori $L^2$ vectors, whereas the induction was stated only for
smooth test functions; an approximation/closability argument was not supplied.  The claimed
implication ``(P-U) $\Rightarrow$ (H6)'' is therefore withdrawn.

\subsection{A corrected conditional recursion}

The mean terms can be summed without $\Tr W_u$, but doing so exposes an additional input.  The
following proposition records the valid statement.

\begin{proposition}[corrected covariance recursion]\label{prop:recursion}
Fix a bounded interval $B$ and an integer $N\ge1$.  Suppose, uniformly in $g$, that:
\begin{enumerate}
\item the covariance-weighted Poincar\'e inequality
\[
 \Var_{\mu_g}(F)\le c_P\E_{\mu_g}\langle\nabla F,C\nabla F\rangle
\]
holds for the ultraviolet-mollified Wick polynomials introduced below;
\item for $0\le n\le N$ the mollified one-point functions satisfy
\[
 m_{n,g,\varepsilon}(x)=
 \E_{\mu_g}\!\left[:\varphi_\varepsilon(x)^n:\right],
 \qquad
 \sup_{\substack{g\in\G\\0<\varepsilon\le1}}
 \|m_{n,g,\varepsilon}\|_{L^\infty(B)}\le b_n;
\]
\item for every $u\in C_c^\infty(B)$ and $n\le N$,
$F_{n,\varepsilon}[u]\to F_n[u]$ in $\mu_g$-probability as $\varepsilon\downarrow0$, where
\[
 F_{n,\varepsilon}[u]=\int_Bu(x):\varphi_\varepsilon(x)^n:\,dx,
 \qquad
 F_n[u]=\int_Bu(x):\varphi(x)^n:\,dx.
\]
\end{enumerate}
Then, for every $n\le N$,
\[
 \sup_{\substack{g\in\G,\ u\in C_c^\infty(B;\R)\\\|u\|_\infty\le1}}
 \E_{\mu_g}[F_n[u]^2]<\infty.
\]
Consequently (M$_2$) holds for every polynomial of degree at most $N$; translation covariance
makes the constant independent of the location of the interval $B$.
\end{proposition}

\begin{proof}
Choose a nonnegative, compactly supported smooth approximate identity $\rho_\varepsilon$ and put
$\varphi_\varepsilon=\rho_\varepsilon*\varphi$ and
\[
 C_\varepsilon=\rho_\varepsilon*C*\widetilde\rho_\varepsilon,
 \qquad \widetilde\rho_\varepsilon(x)=\rho_\varepsilon(-x).
\]
Wick ordering at scale $\varepsilon$ is with respect to $C_\varepsilon(0)$.  Let
$K_{n,g,\varepsilon}$ denote the covariance kernel of the mollified $n$th Wick power, so
\[
 \Var_{\mu_g}(F_{n,\varepsilon}[u])
 =\langle u,K_{n,g,\varepsilon}u\rangle.
\]
All inequalities below are inequalities of positive kernels, equivalently of their quadratic
forms on test functions.  Set
\[
 S_{1,g,\varepsilon}=c_PC_\varepsilon.
\]
For $n\ge2$ define recursively the positive kernel
\[
 S_{n,g,\varepsilon}(x,y)=c_Pn^2C_\varepsilon(x-y)
 \bigl(S_{n-1,g,\varepsilon}(x,y)
 +m_{n-1,g,\varepsilon}(x)m_{n-1,g,\varepsilon}(y)\bigr).       \tag{6.2}
\]
We prove $0\le K_{n,g,\varepsilon}\le S_{n,g,\varepsilon}$ by induction.

For $n=1$, the covariance-weighted norm of the gradient is
$\langle u,C_\varepsilon u\rangle$, so Poincar\'e gives
$\langle u,K_{1,g,\varepsilon}u\rangle
 \le c_P\langle u,C_\varepsilon u\rangle$.
Assume $K_{n-1,g,\varepsilon}\le S_{n-1,g,\varepsilon}$.  The chain rule gives
\[
 \frac{\delta F_{n,\varepsilon}[u]}{\delta\varphi(z)}
 =n\int_Bu(x)\rho_\varepsilon(x-z)
      :\varphi_\varepsilon(x)^{n-1}:\,dx.
\]
Degree lowering and Poincar\'e therefore give
\begin{align*}
 \langle u,K_{n,g,\varepsilon}u\rangle
 &\le c_Pn^2\iint u(x)u(y)C_\varepsilon(x-y)\\
 &\quad\times\bigl(K_{n-1,g,\varepsilon}(x,y)
 +m_{n-1,g,\varepsilon}(x)m_{n-1,g,\varepsilon}(y)\bigr)\,dx\,dy.
\end{align*}
The Schur product theorem applies to positive kernels: multiplying the positive-form inequality
$K_{n-1,g,\varepsilon}\le S_{n-1,g,\varepsilon}$ by the positive-definite kernel
$C_\varepsilon$ preserves the inequality.  The right side is therefore at most
$\langle u,S_{n,g,\varepsilon}u\rangle$, proving the induction.

The ground-state derivative formula is not involved.  Iterating (6.2) instead shows that
$S_{n,g,\varepsilon}$ is a finite sum of terms, each bounded in absolute value on
$B\times B$ by
\[
 A\sum_{r=1}^n C_\varepsilon(x-y)^r,
\]
where $A$ depends only on $n,c_P,b_0,\ldots,b_{n-1}$.  Every finite power of the logarithmically
singular kernel $C(x-y)$ is locally integrable in one dimension, and the bound is uniform in
$\varepsilon$.  Indeed, with
$\kappa_\varepsilon=\rho_\varepsilon*\widetilde\rho_\varepsilon$ a probability density,
$C_\varepsilon=\kappa_\varepsilon*C$; Jensen's inequality gives, for every integer $r\ge1$,
\[
 C_\varepsilon(z)^r\le(\kappa_\varepsilon*C^r)(z).
\]
After enlarging the fixed compact difference set $B-B$ by the support of $\kappa_\varepsilon$,
integration shows
\[
 \sup_{0<\varepsilon\le1}\int_{B-B}C_\varepsilon(z)^r\,dz<\infty.
\]
Hence $\langle |u|,S_{n,g,\varepsilon}|u|\rangle$ is bounded uniformly in $g$ and
$\varepsilon$.  The mean obeys
\[
 \left|\E F_{n,\varepsilon}[u]\right|
 \le b_n\|u\|_{L^1(B)},
\]
so $\sup_{g,\varepsilon}\E[F_{n,\varepsilon}[u]^2]<\infty$.

For fixed $g$, convergence in probability supplies a sequence $\varepsilon_j\downarrow0$ along
which $F_{n,\varepsilon_j}[u]\to F_n[u]$ almost surely.  Fatou's lemma yields
\[
 \E_{\mu_g}[F_n[u]^2]
 \le\liminf_{j\to\infty}\E_{\mu_g}[F_{n,\varepsilon_j}[u]^2],
\]
with the same bound independent of $g$.  A finite sum and Cauchy--Schwarz handle a polynomial
$P$.
\end{proof}

The extra uniform $L^\infty$ hypothesis on the mollified one-point functions in
Proposition~\ref{prop:recursion} is not supplied by
the first-moment estimate quoted in \texttt{h6-attack.tex}.  A bound
$|\E F_n[u]|\le c\|u\|_\infty$ identifies a finite signed measure, not an $L^\infty$ density;
an atomic measure has infinite logarithmic self-energy and would not close (6.2).  This is the
exact distinction hidden by the false inequality (6.1).

There are therefore two honest routes to (M$_2$):
\begin{enumerate}
\item prove a local $N_\tau$/hypercontractive estimate directly, uniform in the remote spatial
cutoff; or
\item prove both the covariance-weighted Poincar\'e inequality and the uniform one-point density
bounds in Proposition~\ref{prop:recursion}.
\end{enumerate}
The unweighted LSI of \cite{BD} supplies neither item 1 nor the covariance-weighted inequality in
item 2.

\section{Corrected proof program}

The dependency order is now:
\[
 \boxed{\text{classical cutoff construction}}
 \Longrightarrow
 \boxed{\text{finite propagation + }\Om_g\in D(V(h))}
\]
and, independently,
\[
 \begin{gathered}
 \boxed{\text{(G$_\mathrm{exact}$)}}
 \Longrightarrow\ \boxed{\text{FKN path regularity, Theorem~\ref{thm:FKNpath}}},\\[5pt]
 \boxed{\text{(G$_\mathrm{exact}$)}}\quad+\quad\boxed{\text{(M$_2$)}}\\[3pt]
 \Longrightarrow\ \boxed{\text{LPPL cell estimate in \texttt{lppl.tex}}}
 \Longrightarrow\ \boxed{\text{full-net local-norm convergence}}.
 \end{gathered}
\]
The last implication includes local normality, translation invariance, $\mathbb Z_2$ invariance
for even $P$, and the algebraic ground-state property, exactly as proved in \texttt{lppl.tex} and
\texttt{../m0/m0.tex}.

The remaining work is to be carried out in the following explicit order.

\begin{enumerate}
\item \textbf{Path theorem (executed).}  Theorem~\ref{thm:FKNpath} proves norm-$C^1$
dependence of $e^{-TK(g_s)}$ directly from the Feynman--Kac--Nelson formula.  Under
(G$_\mathrm{exact}$), spectral perturbation then gives norm-continuous ground-state projections,
the Hellmann--Feynman identity, the ground-state-vector derivative, and endpoint
strong-resolvent continuity.  No common form domain is needed.

\item \textbf{Regime-specific physical gap.}  State the coupling/mass region before proving a
gap.  Establish uniform Euclidean-time connected decay, with one exponent, for a dense cylinder
algebra and for every profile $0\le g\le1$.  Apply Lemma~\ref{lem:decaygap}.  A stochastic-
quantisation spectral gap alone is not a substitute.

\item \textbf{Local second moment.}  Prove (M$_2$) by a uniform local higher-order estimate, or
verify every hypothesis of Proposition~\ref{prop:recursion}.  In the latter route, the required
Poincar\'e form is covariance-weighted and the one-point functions must have uniformly bounded
densities; a total-variation bound is insufficient.

\item \textbf{Assembly (executed).}  The FKN path theorem,
Proposition~\ref{prop:propagation}, and Proposition~\ref{prop:domain} have been inserted into
\texttt{lppl.tex}.  Its cell estimate now assumes only (G$_\mathrm{exact}$) and (M$_2$).
The derivative is integrated on $[\varepsilon,1-\varepsilon]$ and the endpoints are obtained by
norm continuity of the ground-state projection.

\item \textbf{Wavelet uniformity only afterward.}  Once the continuum theorem is unconditional
in a stated regime, prove discrete propagation, gap, and moment estimates with constants uniform
in resolution, and derive path regularity from the corresponding discrete FKN formula.  The
continuum gap and the stochastic-quantisation gap must not be identified.
\end{enumerate}

This ordering removes completed work from the hypothesis list, prevents the endpoint-domain
problem from contaminating the interior derivative, and makes the two genuinely new estimates
--- physical gap and local second moment --- visible as separate theorems.

\section{Final status theorem}

\begin{theorem}[rigorous status of M4]\label{thm:status}
Within the standard spatial-cutoff $P(\varphi)_2$ construction:
\begin{enumerate}
\item the actual cutoff dynamics satisfies exact finite propagation with speed at most one;
\item every cutoff ground state belongs to the domain of every compactly smeared Wick polynomial;
\item (G$_\mathrm{exact}$) implies the affine-path regularity and endpoint continuity needed in
the state-derivative argument;
\item (G$_\mathrm{exact}$) and (M$_2$) imply all conclusions of the conditional continuum LPPL
theorem in \texttt{lppl.tex};
\item neither the uniform physical gap nor (M$_2$) follows from the Bauerschmidt--Dagallier LSI;
\item the proof of ``(P-U) $\Rightarrow$ (H6)'' in \texttt{h6-attack.tex} is invalid for the
reasons in Proposition~\ref{prop:failed}.
\end{enumerate}
Consequently the main spatial-cutoff convergence theorem is not yet unconditional.  Its remaining
inputs are exactly (G$_\mathrm{exact}$) and (M$_2$).
\end{theorem}

\begin{proof}
Items 1 and 2 are Propositions~\ref{prop:propagation} and \ref{prop:domain}.  Item 3 is
Theorem~\ref{thm:FKNpath}.  For item 4, that theorem supplies the derivative formula and endpoint
continuity, Proposition~\ref{prop:propagation} supplies the exact light cone,
(G$_\mathrm{exact}$) supplies the spectral filter and clustering estimate, and (M$_2$) supplies
the uniform cell norm.  These are exactly the inputs used in the proof of the one-path estimate in
\texttt{lppl.tex}; no hypothesis (D) or (L) remains.  Items 5 and 6 were proved in
\S\ref{sec:gap} and Proposition~\ref{prop:failed}.  The last sentence is the resulting exhaustive
list.
\end{proof}

\begin{thebibliography}{9}
\bibitem{BD}
R.~Bauerschmidt and B.~Dagallier,
\emph{Log-Sobolev inequality for the $\varphi^4_2$ and $\varphi^4_3$ measures},
Comm. Pure Appl. Math. \textbf{77} (2024), 2579--2612;
\href{https://arxiv.org/abs/2202.02295}{arXiv:2202.02295}.

\bibitem{GJExp}
J.~Glimm and A.~Jaffe,
\emph{Quantum Field Theory and Statistical Mechanics: Expositions},
Birkh\"auser, Boston, 1985, Chapter~4.

\bibitem{GJModels}
J.~Glimm and A.~Jaffe,
\emph{Boson quantum field models}, in \emph{Mathematics of Contemporary Physics},
Academic Press, 1972, Chapters~7--8; reprinted in \cite{GJExp}.

\bibitem{GJIV}
J.~Glimm and A.~Jaffe,
\emph{The $\lambda(\varphi^4)_2$ quantum field theory without cutoffs. IV.
Perturbations of the Hamiltonian},
J. Math. Phys. \textbf{13} (1972), 1568--1584.

\bibitem{GRS}
F.~Guerra, L.~Rosen, and B.~Simon,
\emph{Nelson's symmetry and the infinite volume behavior of the vacuum in $P(\varphi)_2$},
Comm. Math. Phys. \textbf{27} (1972), 10--22.

\bibitem{Simon}
B.~Simon,
\emph{The $P(\varphi)_2$ Euclidean (Quantum) Field Theory},
Princeton University Press, Princeton, 1974.
\end{thebibliography}

\end{document}
